\documentclass[12pt]{article}
\usepackage[T2A]{fontenc}
\usepackage[utf8]{inputenc}
\usepackage[russian]{babel}
\usepackage{amsmath,amssymb}
\usepackage{geometry}
\geometry{a4paper, margin=2.5cm}

\begin{document}

\section*{Трансформерные нейронные операторы для задачи распространения термоупругих волн}

\subsection*{Постановка}

Требуется построить аппроксимацию оператора решения связанной системы термоупругости
\[
\mathcal{G} : (a, \mathrm{BC}, \mathrm{IC}) \;\longmapsto\; (u, T, \sigma),
\]
где $a(x)$~--- набор параметров среды ($\rho, \lambda, \mu, \gamma, k, c_p$),
$\mathrm{BC}, \mathrm{IC}$~--- граничные и начальные условия,
а $(u, T, \sigma)$~--- предсказываемые поля смещения, температуры и напряжений.
Среда предполагается неоднородной (слоистой) и, в общем случае, анизотропной.

\subsection*{Идея трансформера}

Каждая точка расчётной области представляется в виде \emph{токена}~--- вектора, содержащего её координаты, локальные значения полей и параметры среды:
\[
\mathbf{x}_i = \bigl(\,p_i,\; u_i,\; T_i,\; \rho_i,\; \lambda_i,\; \mu_i,\; \gamma_i,\; k_i,\; c_{p,i}\,\bigr).
\]
Механизм самовнимания (\emph{self-attention}) вычисляет попарные веса взаимодействия между всеми токенами:
\[
\mathrm{Attn}(Q,K,V)
=
\mathrm{softmax}\!\left(\frac{QK^{\top}}{\sqrt{d}}\right) V,
\qquad
Q = X W_Q,\;\; K = X W_K,\;\; V = X W_V,
\]
где $X$~--- матрица токенов, а $W_Q, W_K, W_V$~--- обучаемые матрицы.
В отличие от свёрточных и графовых архитектур, такой механизм обеспечивает \emph{глобальное взаимодействие за один слой}, без накопления через цепочки соседей.

\subsection*{Архитектура для термоупругой задачи}

Архитектура состоит из трёх блоков:
\begin{enumerate}
\item \textbf{Энкодер.} MLP проектирует токены $\mathbf{x}_i$ в латентное пространство размерности $d$.
\item \textbf{Процессор.} Последовательность из $L$ блоков самовнимания с остаточными связями и нормализацией.
\item \textbf{Декодер.} По запросной координате $(x,t)$ вычисляются предсказания
$\hat{u}(x,t)$, $\hat{T}(x,t)$, $\hat{\sigma}(x,t)$.
\end{enumerate}

Функция потерь объединяет ошибку на данных и физические невязки уравнений Куна\,--\,Томпсона:
\[
\mathcal{L}
=
\lambda_{\mathrm{data}}\,\mathcal{L}_{\mathrm{data}}
+
\lambda_{\mathrm{PDE}}\,\mathcal{L}_{\mathrm{PDE}}
+
\lambda_{\mathrm{BC}}\,\mathcal{L}_{\mathrm{BC}}
+
\lambda_{\mathrm{IC}}\,\mathcal{L}_{\mathrm{IC}}.
\]

\subsection*{Преимущества выбора трансформера}

\begin{itemize}
\item \textbf{Глобальное взаимодействие.} Волны распространяются на большие расстояния, и в каждой точке решение зависит от удалённых источников и интерфейсов между слоями. Самовнимание учитывает все эти зависимости одновременно, тогда как у MeshGraphNet радиус взаимодействия ограничен числом шагов передачи сообщений, а у CNN~--- размером рецептивного поля.
\item \textbf{Естественная работа с гетерогенной средой.} Параметры пласта ($\rho, \lambda, \mu, \gamma, k$) подаются непосредственно как признаки токена; модель сама обучается выделять границы слоёв и обрабатывать контрастные интерфейсы без специальной разметки сетки.
\item \textbf{Независимость от регулярной сетки.} Архитектура работает на произвольных множествах точек, что важно для задач с криволинейными границами и адаптивной дискретизацией.
\item \textbf{Связанная многокомпонентная физика.} Перекрёстное внимание между механическими и тепловыми степенями свободы позволяет модели явно выучивать связь
$\gamma T_0\,\partial \varepsilon_{kk}/\partial t$
в уравнении теплопроводности и
$-\gamma \delta_{ij} T$
в определяющем соотношении.
\item \textbf{Высокая скорость инференса.} После однократного обучения предсказание для новой конфигурации среды выполняется за один проход, без необходимости повторного решения уравнений.
\item \textbf{Повторная используемость для обратной задачи.} Та же архитектура без изменений применима к восстановлению параметров среды по наблюдаемым волновым полям.
\end{itemize}

\subsection*{Ограничения}

\begin{itemize}
\item Требуется заранее сгенерированный обучающий набор решений (МКР или МКЭ).
\item Сложность стандартного самовнимания~--- $O(N^2)$ по числу точек; на больших сетках необходимы линейные варианты (Galerkin Transformer, Transolver).
\item Физическая интерпретируемость уступает PINN, где уравнения встроены в функцию потерь.
\end{itemize}

\subsection*{Рекомендуемые варианты}

\begin{itemize}
\item \textbf{OFormer}~--- базовая операторная архитектура; удобна для первичных экспериментов.
\item \textbf{Galerkin Transformer}~--- линейная сложность по числу точек; теоретически обоснован через интегральные операторы.
\item \textbf{Transolver}~--- современный вариант с группировкой точек в физические токены; устойчив на гетерогенных средах с резкими контрастами параметров.
\end{itemize}

\subsection*{Заключение}

Трансформерные нейронные операторы являются обоснованным выбором для задачи прогнозирования распространения термоупругих волн в неоднородных геологических средах. Они обеспечивают глобальное взаимодействие между точками области, естественно работают с гетерогенными параметрами среды и многокомпонентными полями, а также допускают расширение на обратные задачи геофизики.

\end{document}
